% chapters/introduction.tex
\chapter{Introduction}
\label{ch:introduction}

\section{Background and Motivation}
\label{sec:intro-background}

Large Language Models (LLMs) are now a routine presence in academic settings, used
by students and educators to support reasoning, explanation, and decision-making
\citep{miliani2026explica, zhang2023llmeval, lin2024criticbench}. As this use
accelerates, two pressing questions emerge: how do students evaluate the quality of
LLM-generated explanations, and how can different LLMs be compared systematically
across varied reasoning tasks?

Explanation quality varies widely across models, and misleading or incoherent
outputs can erode student trust and distort understanding \citep{kiciman2024causal}.
Yet existing evaluation approaches either treat quality criteria in isolation or
collapse them into a single scalar score \citep{ehsan2020human}, while standard
accuracy-based benchmarks offer little insight into reasoning differences across
task types. This leaves model selection without a principled, human-grounded basis.

The present study addresses both gaps. The first aim is to quantify student judgment
of LLM-generated explanations by identifying which criteria students rely on and how
they weight them. The second aim is to provide a structured, multi-criteria basis for
comparing LLM performance across reasoning task categories, moving beyond scalar
metrics toward a richer characterisation of model strengths and weaknesses.

A key insight motivating this work is that evaluator \emph{disagreement}, typically
discarded as noise, may instead signal model reasoning instability. Where reviewers
disagree most strongly, the underlying model output is likely to be inconsistent or
ambiguous. Preserving this variation, rather than averaging it away, reveals
qualitative differences that accuracy-based metrics miss.

\section{Research Aims and Questions}
\label{sec:intro-aims}

This dissertation is organised around two primary research aims:

\begin{enumerate}[label=\textbf{Aim \arabic*.}, leftmargin=*, labelindent=0pt]
  \item \textbf{Quantify student judgment criteria.} Identify and weight the
    dimensions along which graduate students evaluate the quality of LLM-generated
    causal explanations, using a structured pairwise comparison procedure.
  \item \textbf{Compare LLM performance across reasoning task types.} Provide a
    structured multi-criteria comparison of nine open-source LLMs evaluated on
    economic reasoning tasks, capturing both average performance and variation in
    student ratings.
\end{enumerate}

These aims are operationalised through three concrete research questions:

\begin{enumerate}
  \item Which explanation quality criteria do graduate students prioritise when
    evaluating LLM outputs on economic reasoning tasks?
  \item How do nine open-source LLMs compare in terms of student-rated explanation
    quality across five reasoning task categories?
  \item How does preserving reviewer disagreement through interval-based aggregation
    change the picture relative to simple averaging?
\end{enumerate}

\section{Key Contributions}
\label{sec:intro-contributions}

This work makes four contributions to the evaluation of LLM explanations:

\begin{enumerate}
  \item \textbf{A student-centred evaluation framework} that integrates AHP-derived
    criterion weights with Likert-scale rating data collected from graduate students,
    grounded in the explanation quality dimensions of \citet{hoffman2019metrics}.
  \item \textbf{An adaptive question-selection procedure} based on a composite
    semantic and structural dispersion score that identifies the most discriminative
    items from the MMLU Economics subset, reducing annotation cost without reducing
    coverage.
  \item \textbf{An interval-based aggregation approach} that represents reviewer
    disagreement through interquartile range (IQR) intervals, exposing reasoning
    instability that scalar means conceal.
  \item \textbf{Empirical findings on nine open-source LLMs} showing that
    trustworthiness and completeness dominate student judgment and that a systematic
    fluency trap affects multiple models.
\end{enumerate}

\section{Thesis Structure}
\label{sec:intro-structure}

The remainder of this dissertation is organised as follows. Chapter~\ref{ch:literature}
reviews related work on causal reasoning in LLMs, human-centred explainability, and
multi-criteria decision making. Chapter~\ref{ch:methodology} presents the three-stage
evaluation framework, including data collection, question selection, and aggregation.
Chapter~\ref{ch:results} reports the quantitative and qualitative findings.
Chapter~\ref{ch:discussion} interprets the results and discusses their implications.
Chapter~\ref{ch:conclusion} concludes and outlines directions for future work.
